\chapter{Applicazioni di successo}
Osserviamo ora alcuni modelli di Intelligenza 
Artificiale che si sono dimostrati in grado di produrre risultati concreti e validati 
clinicamente. 
I tre casi documentati di seguito illustrano come i 
metodi trattati in questo elaborato abbiano già cambiato 
la pratica della \textit{drug discovery}.

\section{Baricitinib per COVID-19}
Nel 2020, il team di BenevolentAI ha pubblicato 
su \textit{The Lancet} un'analisi che illustra come una rappresentazione a  
\textbf{knowledge graph} avanzato, costruito come un archivio strutturato di 
informazioni biologiche e mediche estratte tramite ML,
abbia permesso di identificare in pochi giorni il 
\textbf{baricitinib} come potenziale trattamento per 
l'infezione da SARS-CoV-2 
\cite{richardsonBaricitinibPotentialTreatment2020}.

Il baricitinib è un inibitore delle chinasi Janus (JAK), 
già conosciuto e approvato per il trattamento dell'artrite reumatoide. 
Gli algoritmi di predizione DTI di BenevolentAI hanno 
identificato la sua capacità di inibire la proteina 
\textbf{AAK1}, un regolatore dell'endocitosi (il meccanismo attraverso cui il virus entra nelle cellule 
polmonari). Questo ha suggerito un duplice meccanismo 
d'azione: 
\begin{itemize}
    \item La riduzione dell'ingresso virale all'interno dei tessuti polmonari.
    \item La modulazione della risposta infiammatoria nel caso di avvenuta contaminazione.
\end{itemize}
Il farmaco ha successivamente ricevuto l'autorizzazione d'uso 
emergenziale, dimostrando un beneficio significativo 
sulla mortalità nei pazienti con COVID-19 grave.

Questo caso è emblematico del \textbf{drug repurposing 
guidato da ML}, ovvero la capacità di identificare nuovi 
target biologici per farmaci già approvati, 
riducendo drasticamente i tempi di validazione 
clinica.
Si tratta una delle applicazioni più immediate 
e ad efficaci dei metodi di predizione DTI 
discussi in questo elaborato.

\section{Halicin: un antibiotico scoperto dal Deep 
Learning}
Nel 2020, un gruppo di ricercatori del MIT ha pubblicato 
su \textit{Cell} la scoperta di \textbf{Halicin}, il 
primo antibiotico ad ampio spettro identificato tramite 
un approccio di deep learning 
\cite{stokesDeepLearningApproach2020}.

Il modello adottato è \textit{Chemprop}, un'architettura \textbf{D-MPNN} (\textit{Directed Message Passing Neural Network}). Si tratta di una variante delle GNN (descritte nei capitoli precedenti) con lo scopo di prevedere l'attività antibatterica a partire dalla struttura chimica del composto, espressa sotto forma di grafo.

Il modello traduce la rappresentazione in un vettore tramite \textit{message passing}, per costruire una comprensione della chimica locale della molecola.

Il vantaggio rispetto ai metodi tradizionali che fanno uso di fingerprint Morgan o descrittori RDKit (che nel paper vengono usati come \textit{baseline}), è che il \textit{message passing} permette l'apprendimento automatico in funzione della proprietà da predire, senza aver la necessità di introdurre definizioni manuali.
Tuttavia, per aumentare le prestazioni della rete neurale, sono stati concatenate alla rappresentazione iniziale, 200 descrittori chimici classici (da RDKit) per bilanciare la difficile propagazione dei messaggi su grafi molecolari molto estesi.

Successivamente è avvenuta la fase di \textit{training}, svolto per il task di classificazione DTI su molecole in grado di inibire la crescita di \textit{E. coli}. Vengono coinvolte nello screening 2.560 molecole e vengono identificate come \textit{hit} 120 composti attivi (5.14\%).
Dopo deduplicazione, il dataset finale è di 2.335 molecole, binarizzate come hit/non-hit.

Il training viene sottoposto a un \textit{fine tuning} attraverso l'\textbf{ottimizzazione Bayesiana dei parametri}, che usa \textit{Bayesian optimization} con 20 iterazioni, esplorando diverse tipologie di iperparametri.

\begin{table}[ht]
    \centering
    \caption{Elenco degli iperparametri del modello.}
    \label{tab:iperparametri}
    \begin{tabular}{lcc}
        \toprule
        \textbf{Hyperparameter} & \textbf{Range} & \textbf{Value} \\
        \midrule
        Number of message-passing steps & [2, 6] & 5 \\
        Neural network hidden size      & [300, 2400] & 1600 \\
        Number of feed-forward layers   & [1, 3] & 1 \\
        Dropout probability             & [0, 0.4] & 0.35 \\
        \bottomrule
    \end{tabular}
\end{table}

Inoltre, il modello finale utilizza un \textit{ensemble} di 20 istanze pesate in modo diverso per aumentare la robustezza della predizione.

La predizione avviene tramite applicazione sulla \textbf{Drug Repurposing Hub}, completa di 6.111 molecole da cui vengono escluse quelle già presenti nel Training-Set.
Vengono individuate le 99 candidate con punteggi più alti, che passano alla fase di testing sperimentalmente. Di queste, 51 si confermano veri positivi.

Un modello Chemprop specifico viene poi addestrato sul dataset \textbf{ClinTox} per prevedere la tossicità delle molecole candidate rimaste.
Ognuna delle entry del dataset è composta da due proprietà binarie, la tossicità rilevata negli studi clinici e lo stato di approvazione della FDA.

Un \textit{ensemble} di 5 modelli seleziona infine \textbf{Halicin} come antibiotico migliore. Si tratta di un inibitore della c-Jun N-terminal kinase originariamente sviluppato per il diabete.

Il risultato informaticamente più rilevante è che 
Halicin ha una struttura chimica radicalmente diversa 
da tutti gli antibiotici convenzionali. Il modello ha quindi dimostrato capacità di \textbf{scaffold hopping},
apprendendo relazioni struttura-attività che sfuggirebbero
all'intuizione umana dei ricercatori e superando il bias
verso strutture chimiche note tipico dei metodi tradizionali.

Testata su modelli murini, 
Halicin ha eliminato infezioni da 
\textit{Clostridioides difficile} e 
\textit{Acinetobacter baumannii} che resistono 
agli antibiotici convenzionali.

Infine, espandendo lo 
screening al database \textbf{ZINC15}, composto da oltre 107 
milioni di molecole, il modello ha identificato in 4 giorni
ulteriori otto composti antibatterici strutturalmente 
distanti dai farmaci noti \cite{stokesDeepLearningApproach2020}.
\section{Rentosertib: il primo farmaco end-to-end 
generato dall'AI}
Il caso più avanzato e recente è quello del 
\textbf{rentosertib} (ISM001-055), sviluppato da 
Insilico Medicine e documentato in uno studio pubblicato su 
\textit{Nature Medicine} nel giugno 2025 
\cite{xuGenerativeAIDiscoveredTNIK2025}.

A differenza dei casi precedenti, i quali applicano ML 
a composti già esistenti, il rentosertib rappresenta 
il primo farmaco in cui sia il \textbf{target biologico} 
che la \textbf{struttura molecolare} del composto sono stati 
identificati o progettati interamente da algoritmi di 
AI generativa.

La piattaforma \textbf{PandaOmics} ha 
identificato \textbf{TNIK} come target critico per la 
fibrosi polmonare idiopatica (IPF), una patologia 
progressiva e attualmente incurabile. La piattaforma 
di chimica generativa ha poi progettato il suo 
inibitore, completando l'intera pipeline dalla target 
identification al candidato preclinico in 
\textbf{soli 18 mesi} (tempi molto ridotti in confronto alle pipeline tradizionali), con meno di 80 molecole 
sintetizzate e testate.

La generazione dell'inibitore specifico di TNIK è avvenuta tramite la piattaforma di chimica generativa \textit{Chemistry42}. Questa ha messo in parallelo 30 modelli generativi che, basandosi su metriche di valutazione e vincoli strutturali condivisi, hanno permesso un'esplorazione dello spazio chimico collaborativa e priva di ridondanze \cite{renSmallmoleculeTNIKInhibitor2025}.

Il sito di legame selezionato come luogo dell'interazione con il bersaglio è risultato il
\textbf{sito di legame dell'ATP di TNIK}.
La generazione è stata vincolata dalle seguenti ipotesi 
\textbf{farmacoforiche} \cite{renSmallmoleculeTNIKInhibitor2025}: 
\begin{enumerate}
    \item Le strutture 
    candidate devono formare un legame idrogeno con 
    l'NH di \textbf{Cys108} nella regione "hinge" (snodo flessibile) della 
    chinasi.
    \item I centri idrofobici devono essere posizionati nello spazio adiacente al gatekeeper \textbf{Met105} (amminoacido che determina l'accessibilità al sito attivo).
\end{enumerate}
Questi vincoli strutturali, gestiti dai moduli 
\textit{Pocket} e \textit{Pharmacophore} della 
piattaforma, guidano la generazione verso strutture 
con le caratteristiche di binding necessarie 
per l'inibizione dell'ATP.

Infine, algoritmi di screening secondari rifiniscono la selezione scartando le strutture chimicamente instabili o troppo complesse da sintetizzare. Questo stringente processo di filtraggio computazionale converte un immenso spazio di generazione in un set ristretto di molecole ad alto potenziale, massimizzando il tasso di successo della sperimentazione in vitro \cite{renSmallmoleculeTNIKInhibitor2025}.